\begin{derivation}[对称破缺后的径向与角向质量]

正文\eqnrefs{eq:ssb-complex-scalar-lag-dp11} 给出了只依赖 \(|\varphi_{\rm ssb}|\) 的对称势。先由驻点条件确定真空半径，再在所选真空附近展开到二次阶；径向与切向二次系数的差别最终给出正文\eqnrefs{eq:ssb-higgs-goldstone-masses-dp11} 的有质量径向模和无质量切向模。


令 \(r^2\equiv\varphi_{\rm ssb}^*\varphi_{\rm ssb}\)，则
\begin{equation*}
 V(r)=-\mu_{\rm ssb}^2r^2+\frac\lambda{\hbar c}r^4.
\end{equation*}
极值条件为
\begin{align*}
 \frac{\dd V}{\dd r}
 &=-2\mu_{\rm ssb}^2r+\frac{4\lambda}{\hbar c}r^3,\\
 &=2r\left(-\mu_{\rm ssb}^2+\frac{2\lambda}{\hbar c}r^2\right)=0.
\end{align*}
除去不稳定的 \(r=0\)，稳定极小值满足
\begin{equation*}
 r^2=\frac{\mu_{\rm ssb}^2\hbar c}{2\lambda}.
\end{equation*}
定义 \(r=v/\sqrt2\) 即得到
\begin{equation*}
 v^2=\frac{\mu_{\rm ssb}^2\hbar c}{\lambda}.
\end{equation*}

将真空附近参数化代入
\begin{equation*}
 \varphi_{\rm ssb}=\frac1{\sqrt2}(v+h+\ii\pi),
 \qquad
 \varphi_{\rm ssb}^*\varphi_{\rm ssb}=\frac12[(v+h)^2+\pi^2].
\end{equation*}
保留到二次阶：
\begin{align*}
 V
 ={}&-\frac{\mu_{\rm ssb}^2}{2}
 (v^2+2vh+h^2+\pi^2)\\
 &+\frac\lambda{4\hbar c}
 (v^2+2vh+h^2+\pi^2)^2.
\end{align*}
线性 \(h\) 项系数为
\begin{equation*}
 -\mu_{\rm ssb}^2v+\frac{\lambda v^3}{\hbar c}=0,
\end{equation*}
二次项为
\begin{equation*}
 V^{(2)}=\mu_{\rm ssb}^2h^2+0\times\pi^2.
\end{equation*}
与 \(\tfrac12(m^2c^2/\hbar^2)\phi^2\) 比较即得\eqnrefs{eq:ssb-higgs-goldstone-masses-dp11}。
\end{derivation}

\begin{derivation}[规范固定怎样消去 Goldstone 混合项]
直角坐标展开给出
\[
\Lag_{\rm mix}=g_{\rm G}v\,\mathscr C^\mu\partial_\mu\pi,
\qquad
F=\partial\!\cdot\!\mathscr C-\xi\frac{g_{\rm G}v}{\hbar c}\pi .
\]
规范固定项展开为
\[
\Lag_{\rm gf}
=-\frac{\hbar c}{2\xi}(\partial\!\cdot\!\mathscr C)^2
+g_{\rm G}v\,\pi\,\partial\!\cdot\!\mathscr C
-\frac{\xi g_{\rm G}^2v^2}{2\hbar c}\pi^2 .
\]
对中间项分部积分，
\[
g_{\rm G}v\,\pi\,\partial_\mu\mathscr C^\mu
=-g_{\rm G}v\,\mathscr C^\mu\partial_\mu\pi+\text{边界项},
\]
故它与 $\Lag_{\rm mix}$ 相消。剩余 $\pi^2$ 项与标准质量项比较：
\[
-\frac{\xi g_{\rm G}^2v^2}{2\hbar c}\pi^2
=-\frac12\left(\frac{m_\pi c}{\hbar}\right)^2\pi^2,
\qquad
\left(\frac{m_Ac}{\hbar}\right)^2=\frac{g_{\rm G}^2v^2}{\hbar c},
\]
于是
\[
\boxed{m_\pi^2=\xi m_A^2}.
\]
$\pi$ 的质量依赖规范参数 $\xi$，而物理规范玻色子质量 $m_A$ 不随 $\xi$ 改变。
\end{derivation}

\begin{derivation}[标量四次相互作用的一圈泡积分]
\eqnrefs{eq:dim-reg-measure-general-dp10} 中显式 \(\hbar\) 的消去可由变量替换直接看出。若从逆长度圈变量 \(\kappa^\mu\) 开始，定义物理动量

\begin{equation*}
 \ell^\mu=\hbar\kappa^\mu,
 \qquad
 \mu_{\rm DR}=\hbar\mu_\kappa,
\end{equation*}
则
\begin{equation*}
 \dd^d \kappa=\frac{\dd^d \ell}{\hbar^d},
 \qquad
 (\kappa^2-M^2)^{-2}
 =\frac{\hbar^4}{(\ell^2-m_R^2c^2)^2},
\end{equation*}
其中 \(M=m_Rc/\hbar\)。维数正则化尺度给
\begin{equation*}
 \mu_\kappa^{2\epsilon}
 =\frac{\mu_{\rm DR}^{2\epsilon}}{\hbar^{2\epsilon}}.
\end{equation*}
由于 \(d=4-2\epsilon\)，总的 \(\hbar\) 次数为
\begin{equation*}
 4-d-2\epsilon=0.
\end{equation*}
因此变量换元后的物理四动量形式为
\begin{equation*}
 \mu_{\rm DR}^{2\epsilon}
 \int\frac{\dd^d \ell}{(2\pi)^d}
 \frac{\ii}{(\ell^2-m_R^2c^2+\ii0^+)
 ((\ell+p)^2-m_R^2c^2+\ii0^+)},
\end{equation*}
其中显式 $\hbar$ 已由测度、分母和正则化尺度的幂次精确抵消。

用 Feynman 参数
\begin{equation*}
 \frac1{AB}=\int_0^1\frac{\dd x}{[xA+(1-x)B]^2}
\end{equation*}
并令 \(L=\ell+xp\)，分母化为
\begin{equation*}
 L^2-\Delta_s(x)+\ii0^+,
 \qquad
 \Delta_s(x)=m_R^2c^2-x(1-x)s.
\end{equation*}
于是
\begin{equation*}
 I(s)=\int_0^1\dd x\,
 \mu_{\rm DR}^{2\epsilon}
 \int\frac{\dd^d L}{(2\pi)^d}
 \frac{\ii}{(L^2-\Delta_s+\ii0^+)^2}.
\end{equation*}
做 Wick 转动 \(L^0=\ii L_E^0\)：
\begin{equation*}
 \dd^d L=\ii\dd^d L_E,
 \qquad
 L^2-\Delta+\ii0^+=-(L_E^2+\Delta-\ii0^+),
\end{equation*}
故
\begin{equation*}
 I(s)=-\int_0^1\dd x\,
 \mu_{\rm DR}^{2\epsilon}
 \int\frac{\dd^d L_E}{(2\pi)^d}
 \frac1{(L_E^2+\Delta_s-\ii0^+)^2}.
\end{equation*}
利用欧氏母积分
\begin{equation*}
 \int\frac{\dd^d L_E}{(2\pi)^d}
 \frac1{(L_E^2+\Delta)^2}
 =\frac1{(4\pi)^{d/2}}
 \Gamma\!\left(2-\frac d2\right)
 \Delta^{d/2-2},
\end{equation*}
代入 \(d=4-2\epsilon\)：
\begin{equation*}
 I(s)=-\int_0^1\dd x\,
 \frac1{(4\pi)^2}
 \Gamma(\epsilon)(4\pi)^\epsilon
 \left(\frac{\mu_{\rm DR}^2}{\Delta_s-\ii0^+}\right)^\epsilon.
\end{equation*}
逐项展开
\begin{align*}
 \Gamma(\epsilon)&=\frac1\epsilon-\gamma_E+\order(\epsilon),\\
 (4\pi)^\epsilon&=1+\epsilon\ln4\pi+\order(\epsilon^2),\\
 \left(\frac{\mu_{\rm DR}^2}{\Delta_s-\ii0^+}\right)^\epsilon
 &=1+\epsilon\ln\frac{\mu_{\rm DR}^2}{\Delta_s-\ii0^+}
 +\order(\epsilon^2).
\end{align*}
相乘到有限阶并执行参数积分，得到
\begin{equation*}
 I(s)=-\frac1{16\pi^2}
 \left[
 \frac1{\bar\epsilon}
 -\int_0^1\dd x\,
 \ln\frac{\Delta_s(x)-\ii0^+}{\mu_{\rm DR}^2}
 \right]+
 \order(\epsilon),
\end{equation*}
即\eqnrefs{eq:phi4-bubble-result-dp10}。当 \(\Delta_s(x)\) 穿过零时，\(-\ii0^+\) 决定对数的虚部，并与在壳中间态的 Cutkosky 不连续度一致。
\end{derivation}

\begin{derivation}[标量四次耦合的一圈演化]
记

\begin{equation*}
 a_1\equiv\frac{3}{32\pi^2},
 \qquad
 P_\epsilon\equiv\frac1{\bar\epsilon}.
\end{equation*}
一圈阶裸耦合为
\begin{equation*}
 g_{4,B}
 =\mu^{2\epsilon}
 \left[g_4(\mu)+a_1P_\epsilon g_4^2(\mu)
 +\order(g_4^3)\right].
\end{equation*}
按正文\eqnrefs{eq:phi4-beta-one-loop-dp10} 中 $\beta_{g_4}$ 的定义，$d=4-2\epsilon$ 时固定裸耦合的总尺度导数写成
\begin{equation*}
 \mu\frac{\dd g_4}{\dd\mu}\Big|_{g_{4,B}}
 =-2\epsilon g_4+\beta_{g_4}(g_4).
\end{equation*}
对裸表达式求导时，反项同样通过 \(g_4(\mu)\) 依赖于 \(\mu\)：
\begin{align*}
 0={}&2\epsilon
 \left(g_4+a_1P_\epsilon g_4^2\right)\\
 &+\left(1+2a_1P_\epsilon g_4\right)
 \left(-2\epsilon g_4+\beta_{g_4}\right)
 +\order(g_4^3).
\end{align*}
保留到 \(g_4^2\) 后，\(2\epsilon g_4\) 项相消，并得到
\begin{equation*}
 0=\beta_{g_4}
 -2a_1\epsilon P_\epsilon g_4^2
 +\order(g_4^3,\epsilon g_4^2).
\end{equation*}
因为
\begin{equation*}
 \epsilon P_\epsilon
 =\epsilon\left(\frac1\epsilon-\gamma_E+\ln4\pi\right)
 \longrightarrow1,
\end{equation*}
四维极限给出
\begin{equation*}
 \beta_{g_4}=2a_1g_4^2+
\order(g_4^3)
 =\frac{3g_4^2}{16\pi^2}+
\order(g_4^3).
\end{equation*}
这里关键的因子 2 来自 \(d=4-2\epsilon\) 下的 \(\mu^{2\epsilon}\)；若把反项对 \(g_4(\mu)\) 的隐含尺度依赖漏掉，就会错误地得到相反符号。
\end{derivation}

\begin{derivation}[固有时积分中的一圈紫外发散]

对正定欧几里得二次算符 $\Delta$，一圈实玻色 高斯行列式可用正文已经建立的固有时表示写成
\[
\frac12\Tr\ln\Delta
=-\frac12\mu_{\rm DR}^{2\epsilon}
\int_0^\infty\frac{\dd s}{s}\,
\Tr \ee^{-s\Delta},
\]
其中与场无关的加法常数不影响局域反项。把正文\eqnrefs{eq:heat-kernel-a4-flat-dp11} 代入，并只保留含 $s^2$ 的四维局域系数，得到
\begin{align*}
\left.\frac12\Tr\ln\Delta\right|_{a_4}
={}&-\frac12\mu_{\rm DR}^{2\epsilon}
\int\dd^d x\,
\tr a_4(x)\\
&\times\int_0^\infty\frac{\dd s}{s}
\frac{s^2}{(4\pi s)^{d/2}}.
\end{align*}
取 $d=4-2\epsilon$ 后，短固有时端的幂次变为
\[
\frac1s\frac{s^2}{s^{2-\epsilon}}
=s^{-1+\epsilon}.
\]
因此紫外极点只需从 $s\to0^+$ 端读取。任取固定分界尺度 $s_0$，
\begin{align*}
\mu_{\rm DR}^{2\epsilon}
\int_0^{s_0}\dd s\,s^{-1+\epsilon}
&=\mu_{\rm DR}^{2\epsilon}
\frac{s_0^\epsilon}{\epsilon}\\
&=\frac1\epsilon
+\ln(\mu_{\rm DR}^2s_0)+O(\epsilon).
\end{align*}
所以极点系数与任意分界 $s_0$ 无关，而长固有时部分只改变有限项。代入
\[
 a_4(x)
 =\frac12E^2
 +\frac1{12}\Omega_{\mu\nu}\Omega^{\mu\nu}
 +\frac16\nabla^2E
\]
并使用无边界体积分
\(
\int\dd^4 x\,\nabla^2E=0
\)，得到
\[
\left.\frac12\Tr\ln\Delta\right|_{\rm pole}
=-\frac1{2(4\pi)^2}\frac1\epsilon
\int\dd^4 x\,
\tr\left[
\frac12E^2+\frac1{12}\Omega_{\mu\nu}\Omega^{\mu\nu}
\right].
\]
最后使用正文已经固定的 $1/\bar\epsilon$ 记号整理 $1/\epsilon-\gamma_E+\ln4\pi$，即得到\eqnrefs{eq:heat-kernel-divergence-master-dp11}。
\end{derivation}

\begin{derivation}[Abel Higgs 模型中的矢量质量]
从
\[
\varphi_{\rm ssb}(x)
=\frac{v+h(x)}{\sqrt2}\ee^{\ii\pi(x)/v},
\qquad
D_\mu=\partial_\mu+\ii g_{\rm G}\mathscr C_\mu
\]
出发。局域 $U(1)$ 变换
\[
\varphi_{\rm ssb}\to \ee^{\ii\alpha(x)}\varphi_{\rm ssb},
\qquad
\mathscr C_\mu\to
\mathscr C_\mu-\frac1{g_{\rm G}}\partial_\mu\alpha
\]
要求极坐标变量变成
\[
h\to h,
\qquad
\pi\to\pi+v\alpha.
\]
因此组合
\[
\mathcal B_\mu
\equiv
\mathscr C_\mu+\frac1{g_{\rm G}v}\partial_\mu\pi
\]
满足 $\mathcal B_\mu\to\mathcal B_\mu$，它才是破缺相位中的规范不变矢量变量。

显式求协变导数：
\begin{align*}
D_\mu\varphi_{\rm ssb}
&=\frac{\ee^{\ii\pi/v}}{\sqrt2}
\left[
\partial_\mu h
+\ii(v+h)
\left(\frac{\partial_\mu\pi}{v}+g_{\rm G}\mathscr C_\mu\right)
\right]\\
&=\frac{\ee^{\ii\pi/v}}{\sqrt2}
\left[
\partial_\mu h
+\ii g_{\rm G}(v+h)\mathcal B_\mu
\right].
\end{align*}
取模平方时，实的径向项与纯虚的相位项交叉项严格相消：
\[
(D_\mu\varphi)^*(D^\mu\varphi)
=\frac12(\partial h)^2
+\frac{g_{\rm G}^2}{2}(v+h)^2\mathcal B_\mu\mathcal B^\mu.
\]
展开 $(v+h)^2=v^2+2vh+h^2$。与矢量二次型有关的纯二次项是
\[
\frac{g_{\rm G}^2v^2}{2}\mathcal B_\mu\mathcal B^\mu.
\]
在上述国际单位制归一化下，Proca 质量项写成
\[
\frac{\hbar c}{2}
\left(\frac{m_Ac}{\hbar}\right)^2
\mathcal B_\mu\mathcal B^\mu.
\]
两者是同一个二次算符，所以系数必须相等：
\[
g_{\rm G}^2v^2
=\hbar c\left(\frac{m_Ac}{\hbar}\right)^2,
\]
从而
\[
\boxed{m_Ac^2=g_{\rm G}\sqrt{\hbar c}\,v.}
\]
Goldstone 场 $\pi$ 通过 $\mathcal B_\mu$ 重组为有质量矢量场的纵向偏振。破缺前自由度计数为“复标量的两个实自由度 + 无质量矢量的两个横向偏振”共四个；破缺后为“径向标量 $h$ 一个 + 有质量矢量 $\mathcal B_\mu$ 三个偏振”仍共四个。这个计数是上述变量重组没有丢失物理自由度的独立检查。
\end{derivation}
